% Laget automatisk av elevark/lag_elevark.py fra sim/skraaplan.html. Endringer her blir overskrevet.
\documentclass[11pt]{article}
\usepackage{elevark}
\begin{document}
\elevarkhode{Fysikk 2 · Bevegelse og krefter i to dimensjoner}{Skråplan med friksjon}{https://fysikklab.no/sim/skraaplan.html}{Simuleringen}

En kloss ligger på et skråplan. Tyngden kan deles i en komponent langs planet, som vil dra klossen nedover, og en komponent normalt på planet, som bestemmer normalkraften og dermed friksjonen. Når begynner klossen å gli, og hvor fort akselererer den?

\section*{Oppgaver}

Start med tre mot tre og 2,0~kg.

\begin{oppgaveliste}
\oppgave \fase{gjett}{Gjett}Ved hvilken helning tror du klossen begynner å gli? Øk vinkelen gradvis og finn ut.
\svar
\oppgave \fase{forklar}{Forklar}Vis at klossen begynner å gli når $\tan  θ = μ_{\mathrm{s}}$. Hvorfor har ikke massen noe å si?
\svar
\oppgave \fase{test}{Regn}Sett $θ = 30^{\circ}$. Regn ut $G \sin  θ$, normalkraften og glidefriksjonen. Vis at $a = g(\sin  θ - μ_{\mathrm{k}} \cos  θ)$, og regn ut akselerasjonen. Sjekk.
\svar
\oppgave \fase{test}{Undersøk}Gi klossen en startfart på 3,0~m/s oppover ved 30°. Hvorfor bremser den raskere på vei opp enn den akselererer på vei ned? Regn ut hvor langt opp den kommer.
\svar
\oppgave \fase{test}{Regn}Hvor stor fart har klossen når den er tilbake der den startet? Sammenlign med startfarten, og forklar hvor energien har blitt av.
\svar
\oppgave \fase{test}{Undersøk}Sett $θ = 20^{\circ}$ og startfart 3,0~m/s. Hva skjer når klossen har stoppet? Hvorfor?
\svar
\oppgave \fase{utfordring}{Utfordring}Uten friksjon: Hvor stor helning gir en akselerasjon på $g/2$? Hvorfor kan akselerasjonen aldri bli større enn $g$?
\svar
\end{oppgaveliste}

\end{document}
